\documentclass[12pt,a4paper]{article}

% ── Packages ──────────────────────────────────────────────────────────────────
\usepackage[a4paper, margin=2.5cm, top=3cm, bottom=3cm]{geometry}
\usepackage[T1]{fontenc}
\usepackage[utf8]{inputenc}
\usepackage{xcolor}
\usepackage{booktabs}
\usepackage{array}
\usepackage{colortbl}
\usepackage{enumitem}
\usepackage{parskip}
\usepackage{titlesec}
\usepackage{fancyhdr}
\usepackage{hyperref}
\usepackage{tcolorbox}
\usepackage{caption}
\usepackage[nottoc]{tocbibind}

% ── Colours ───────────────────────────────────────────────────────────────────
\definecolor{darkblue}{HTML}{1F4E79}
\definecolor{medblue}{HTML}{2E74B5}
\definecolor{lightblue}{HTML}{D6E4F0}
\definecolor{altrow}{HTML}{EBF3FB}
\definecolor{greenbg}{HTML}{E2EFDA}
\definecolor{headerbg}{HTML}{1F4E79}

% ── tcolorbox ─────────────────────────────────────────────────────────────────
\tcbuselibrary{skins, breakable}

\newtcolorbox{keybox}[1]{
  colback=greenbg, colframe=darkblue, coltitle=white,
  fonttitle=\bfseries, title={#1},
  breakable, boxrule=0.8pt, arc=3pt,
  left=6pt, right=6pt, top=4pt, bottom=4pt
}
\newtcolorbox{notebox}{
  colback=lightblue!60, colframe=medblue,
  boxrule=0.6pt, arc=2pt,
  left=6pt, right=6pt, top=3pt, bottom=3pt
}

% ── Section formatting ────────────────────────────────────────────────────────
\titleformat{\section}{\Large\bfseries\color{darkblue}}{\thesection.}{0.6em}{}
  [\color{medblue}\titlerule]
\titleformat{\subsection}{\large\bfseries\color{medblue}}{\thesubsection}{0.5em}{}
\titleformat{\subsubsection}{\normalsize\bfseries\color{darkblue!80}}
  {\thesubsubsection}{0.5em}{}
\titlespacing*{\section}{0pt}{14pt}{8pt}
\titlespacing*{\subsection}{0pt}{10pt}{5pt}
\titlespacing*{\subsubsection}{0pt}{8pt}{4pt}

% ── Header / Footer ───────────────────────────────────────────────────────────
\pagestyle{fancy}
\fancyhf{}
\renewcommand{\headrulewidth}{0.4pt}
\renewcommand{\footrulewidth}{0.4pt}
\fancyhead[L]{\small\color{darkblue}\textbf{Literature Review \& Benchmark Survey}}
\fancyhead[R]{\small\color{darkblue}Rajib Roy --- 24459 $|$ IISc M.Tech (AI)}
\fancyfoot[C]{\small\thepage}
\fancyfoot[R]{\small\color{gray}Phase 1 Deliverable 1 --- April 2026}

% ── Hyperref ──────────────────────────────────────────────────────────────────
\hypersetup{colorlinks=true, linkcolor=medblue, urlcolor=medblue,
  pdftitle={Literature Review and Benchmark Survey --- Rajib Roy},
  pdfauthor={Rajib Roy}}

% ── Table column types ────────────────────────────────────────────────────────
\newcolumntype{L}[1]{>{\raggedright\arraybackslash}p{#1}}
\newcolumntype{C}[1]{>{\centering\arraybackslash}p{#1}}

% ── Lists ─────────────────────────────────────────────────────────────────────
\setlist[itemize]{leftmargin=1.4em, itemsep=3pt, topsep=4pt}
\setlist[enumerate]{leftmargin=1.4em, itemsep=3pt, topsep=4pt}

% ═════════════════════════════════════════════════════════════════════════════
\begin{document}
\pagenumbering{roman}

% ── Title Page ────────────────────────────────────────────────────────────────
\begin{titlepage}
  \thispagestyle{empty}
  \begin{center}
    \vspace*{1cm}
    {\large\color{darkblue}M.Tech.\ (Online) --- Artificial Intelligence}\\[4pt]
    {\large\color{darkblue}Indian Institute of Science, Bangalore}\\[2.5cm]
    {\Huge\bfseries\color{darkblue}Literature Review \&}\\[6pt]
    {\Huge\bfseries\color{darkblue}Benchmark Survey}\\[1.2cm]
    {\color{medblue}\rule{\linewidth}{2pt}}\\[0.6cm]
    {\LARGE\bfseries\color{medblue}
      Detection of AI-Generated and Cloned Speech\\[6pt]
      to Strengthen Voice-Based Authentication\\[6pt]
      in Banking Environments}\\[0.4cm]
    {\color{medblue}\rule{\linewidth}{1pt}}\\[2cm]
    \begin{tabular}{@{}L{4.2cm}L{9cm}@{}}
      \rowcolor{lightblue}\textbf{Student Name}   & Rajib Roy \\
      \rowcolor{white}    \textbf{Student Number} & 24459 \\
      \rowcolor{lightblue}\textbf{Programme}      & M.Tech.\ in Artificial Intelligence \\
      \rowcolor{white}    \textbf{Internal Guide} & Dr.\ Anil Rahate, Practice Head --- AI, Wipro Ltd \\
      \rowcolor{lightblue}\textbf{IISc Mentor}    & Prof.\ Chandra Sekhar Seelamantula,
                                                     Dept.\ of Electrical Engineering, IISc \\
      \rowcolor{white}    \textbf{Deliverable}    & Phase 1 --- Deliverable 1 \\
      \rowcolor{lightblue}\textbf{Date}           & April 2026 \\
    \end{tabular}
  \end{center}
\end{titlepage}

\tableofcontents
\clearpage
\pagenumbering{arabic}

% ═════════════════════════════════════════════════════════════════════════════
\section{Overview and Scope}
% ═════════════════════════════════════════════════════════════════════════════

This document is \textbf{Phase 1 Deliverable 1} of the M.Tech.\ project. It provides a
structured survey of the published literature on synthetic speech detection (commonly
framed as ``anti-spoofing'' in the automatic speaker verification community), with a focus
on methods, datasets, evaluation protocols, and robustness findings relevant to the
project's banking telephony deployment context.

\textbf{Scope of this document:} This is a survey of \emph{existing published work}. It
does not report results from this project's own models, dataset, or experiments, as those
activities are planned for Phase 2. The purpose is to establish what the published literature
says, identify gaps, and record the design decisions taken as a consequence --- decisions
that will govern the methodology implemented in Phases 2 and 3.

The survey covers five areas:
\begin{enumerate}
  \item The ASVspoof challenge series --- the field's primary benchmarking platform.
  \item Classical and handcrafted feature approaches.
  \item Deep learning approaches.
  \item Robustness, codec effects, and domain mismatch.
  \item Gaps in the existing literature and how this project addresses them.
\end{enumerate}

% ═════════════════════════════════════════════════════════════════════════════
\section{The ASVspoof Challenge Series}
% ═════════════════════════════════════════════════════════════════════════════

\subsection{Background}

The ASVspoof challenge series, established in 2015 and running through 2021, is the
primary benchmarking platform for synthetic speech detection (Nautsch et al., 2021). Each
edition defined standardised datasets, evaluation protocols, and baseline systems, enabling
reproducible comparison across research groups.

The challenge distinguishes two principal attack scenarios:

\begin{itemize}
  \item \textbf{Logical Access (LA):} Attacks generated by text-to-speech (TTS) synthesis
    or voice conversion algorithms. This is the attack type most relevant to banking fraud.
  \item \textbf{Physical Access (PA):} Replay attacks, where genuine recordings are
    replayed through a loudspeaker at a microphone.
\end{itemize}

\subsection{ASVspoof 2019}

ASVspoof 2019 (Todisco et al., 2019) is the most directly relevant edition to this
project. Its LA dataset covers 19 TTS and voice-conversion systems (A01--A19), providing
broad coverage of synthesis architectures from that period.

Key contributions:
\begin{itemize}
  \item The \textbf{Constant-Q Cepstral Coefficient (CQCC)} feature was introduced as the
    official baseline, achieving EER $\approx 0.10$ with a GMM classifier on the LA track.
    CQCC provides finer frequency resolution at low frequencies than standard MFCCs, better
    capturing fine spectral structure in vocoder output.
  \item The \textbf{Equal Error Rate (EER)} --- the point on the detection error tradeoff
    curve where the false acceptance rate equals the false rejection rate --- was established
    as the primary evaluation metric for the field.
  \item ASVspoof 2019 generated the most widely-used benchmark dataset for subsequent
    anti-spoofing research, covering statistical parametric, neural, and hybrid TTS systems.
\end{itemize}

\subsection{ASVspoof 2021}

ASVspoof 2021 extended the challenge to evaluate systems under \textbf{matched and
mismatched channel conditions}:

\begin{itemize}
  \item Systems trained on clean studio recordings degraded substantially under telephone
    channel evaluation, with EER increases of 5--15 percentage points reported across
    participating systems.
  \item The challenge demonstrated that codec compression and channel noise modify both the
    genuine speech characteristics and the spoofing artefacts the detectors rely upon,
    making clean-data-trained systems unreliable in real deployments.
  \item This finding represents a critical gap between published benchmark performance and
    real-world banking telephony deployment, where audio is transmitted over G.711 or
    similar codecs at 8\,kHz.
\end{itemize}

\subsection{Detection Cost Function (t-DCF)}

Kinnunen et al.\ (2020) introduced the tandem Detection Cost Function (t-DCF) as a more
operationally realistic performance measure:

\begin{itemize}
  \item Unlike EER, t-DCF jointly assesses the combined error rate of a spoofing
    countermeasure and an ASV system, accounting for the prior probability of attack and
    the relative cost of false acceptance versus false rejection.
  \item This reflects the fact that a countermeasure is not used in isolation --- it
    operates as a front-end to an existing speaker verification system.
  \item t-DCF has been adopted as a co-primary metric alongside EER in ASVspoof 2021.
\end{itemize}

% ═════════════════════════════════════════════════════════════════════════════
\section{Classical and Handcrafted Feature Approaches}
% ═════════════════════════════════════════════════════════════════════════════

\subsection{Systematic Feature Comparison}

Sahidullah et al.\ (2015) provided the most comprehensive comparison of acoustic features
for anti-spoofing at the time of ASVspoof 2015:

\begin{itemize}
  \item Features capturing \textbf{fine spectral structure} consistently outperformed
    standard 13-coefficient MFCCs against vocoder-based attacks. Specifically, using 40
    MFCC coefficients (capturing higher-order cepstral details) outperformed the
    conventional 13-coefficient setting by a substantial margin.
  \item The inclusion of \textbf{delta and delta-delta coefficients} (encoding temporal
    rate-of-change of the spectral envelope) provided consistent gains across all tested
    classifiers. This is attributed to the difference in spectral dynamics between genuine
    articulation --- which produces continuously evolving spectral envelopes --- and TTS
    vocoders, which impose smoother and more regular temporal evolution.
  \item \textbf{Harmonic-to-Noise Ratio (HNR)} and \textbf{spectral shape features}
    (centroid, bandwidth, rolloff) were identified as informative secondary features,
    reflecting the characteristically clean harmonic structure of TTS vocoders.
\end{itemize}

\subsection{Phase Spectrum Features}

Patel and Patil (2015) demonstrated that TTS vocoders introduce \textbf{phase coherence
artefacts} that are invisible in the magnitude spectrum:

\begin{itemize}
  \item Vocoders generate speech with artificially coherent phase structures, whereas
    natural speech exhibits complex and partially random phase relationships across frames.
  \item Features derived from the instantaneous phase spectrum, group delay, or phase
    difference between adjacent frames were shown to capture vocoder artefacts that
    MFCC-based features miss entirely.
  \item Combining phase-based features with MFCC features yielded additive accuracy
    improvements on ASVspoof 2015, motivating their inclusion in broader feature sets.
\end{itemize}

\subsection{SVM Classifiers for Anti-Spoofing}

Support Vector Machines with RBF kernels have been the dominant classifiers for
handcrafted features throughout the ASVspoof 2015--2019 era:

\begin{itemize}
  \item SVMs generalise well with relatively small training sets, compared to the large
    datasets needed by deep learning models, and have strong theoretical foundations for
    binary classification in high-dimensional feature spaces.
  \item The RBF kernel can model non-linear decision boundaries without requiring explicit
    feature interaction terms.
  \item Probabilistic calibration via Platt scaling enables threshold tuning at the EER
    operating point, which is important for class-imbalanced detection problems.
  \item The official ASVspoof 2019 baseline used SVM with CQCC features, and SVM remains
    a strong and reproducible reference system in the field.
\end{itemize}

\subsection{Summary of Findings}

\begin{keybox}{Key Takeaways --- Classical Features}
  \begin{itemize}
    \item Delta and delta-delta MFCCs are the strongest discriminators: they capture
      temporal rate-of-change, which is smoother and more regular in TTS output.
    \item 40 MFCC coefficients substantially outperform the conventional 13-coefficient
      setting for vocoder detection.
    \item Phase coherence features capture artefacts invisible in the magnitude spectrum,
      providing complementary and non-redundant discriminative information.
    \item SVM with RBF kernel is a well-validated, reproducible classical baseline.
  \end{itemize}
\end{keybox}

% ═════════════════════════════════════════════════════════════════════════════
\section{Deep Learning Approaches}
% ═════════════════════════════════════════════════════════════════════════════

\subsection{CNN-Based Systems}

The transition to deep learning in anti-spoofing accelerated with ASVspoof 2019.
Alzantot et al.\ (2019) showed that a ResNet architecture trained end-to-end on log-Mel
spectrograms could match the CQCC + SVM official baseline:

\begin{itemize}
  \item The ResNet learns hierarchical spectral representations directly from the input
    spectrogram, eliminating hand-engineered features.
  \item The model used \textbf{global average pooling across the time dimension} to
    produce a fixed-length representation from variable-length utterances.
  \item This temporal pooling operation discards the order and dynamics of the spectral
    sequence, retaining only average spectral content. This is a known limitation when
    the discriminative signal includes temporal patterns --- which vocoder artefacts do.
\end{itemize}

\subsection{Hybrid CNN-LSTM Architectures}

Multiple subsequent studies have found that hybrid CNN-LSTM architectures consistently
outperform purely convolutional models for anti-spoofing:

\begin{itemize}
  \item The \textbf{CNN front-end} acts as a local spectral feature extractor, identifying
    frequency-domain artefacts within short analysis windows.
  \item The \textbf{LSTM layer} then models the temporal evolution of these spectral
    features across the utterance, capturing prosodic regularity and temporal dynamics
    that CNN global pooling discards.
  \item Using \textbf{frequency-only pooling} (pooling across frequency but not time)
    after the CNN preserves the temporal sequence, enabling the LSTM to process the
    full temporal context.
  \item Bidirectional LSTMs, which process the sequence in both forward and backward
    directions, capture dependencies that span the full utterance and have shown
    further improvements.
\end{itemize}

\subsection{Self-Supervised Representations: wav2vec 2.0}

Baevski et al.\ (2020) introduced wav2vec 2.0, a transformer-based model pre-trained
on 960 hours of unlabelled speech via a contrastive self-supervised objective:

\begin{itemize}
  \item When the frozen or fine-tuned wav2vec 2.0 encoder is used for spoofing detection,
    it achieves competitive EERs without any task-specific feature engineering, suggesting
    that self-supervised pre-training learns representations that are also discriminative
    for synthetic speech detection.
  \item Full fine-tuning requires substantial GPU memory (16\,GB+) and training time.
    Feature extraction from a frozen encoder is computationally tractable and still
    competitive.
  \item This line of work suggests that large pre-trained speech models encode
    properties of natural speech that synthetic systems fail to fully replicate.
\end{itemize}

\subsection{Graph Attention: AASIST}

Wang et al.\ (2021) introduced AASIST (Audio Anti-Spoofing using Integrated
Spectro-Temporal Graph Attention Networks):

\begin{itemize}
  \item AASIST applies graph attention networks to model relationships between nodes
    spanning both spectral and temporal dimensions simultaneously, capturing
    cross-frequency and cross-time dependencies not modelled by CNNs or LSTMs
    independently.
  \item It achieves EER $\approx 0.02$ on ASVspoof 2021 LA --- the state of the art
    at the time of publication.
  \item With approximately 300,000 parameters, AASIST is relatively compact for its
    performance level, making it practically deployable.
\end{itemize}

\subsection{Comparative Published Performance}

Table~\ref{tab:lit_comparison} summarises indicative EER values from key published
anti-spoofing systems on ASVspoof datasets. These numbers represent the state of the
art surveyed during Phase 1 and provide context for understanding the field's progress.

\medskip
\begin{center}
\begin{tabular}{L{3.2cm}C{2.8cm}L{3.8cm}C{1.6cm}L{2.6cm}}
  \toprule
  \rowcolor{headerbg}
  \color{white}\textbf{Reference} &
  \color{white}\textbf{Dataset} &
  \color{white}\textbf{Method} &
  \color{white}\textbf{EER $\downarrow$} &
  \color{white}\textbf{Notes} \\
  \midrule
  Todisco et al.\ (2019) &
  ASVspoof 2019 LA & CQCC + GMM & $\sim$0.10 & Official baseline \\
  \rowcolor{altrow}
  Alzantot et al.\ (2019) &
  ASVspoof 2019 LA & ResNet spectrogram & $\sim$0.05 & CNN-based \\
  Wang et al.\ (2021) &
  ASVspoof 2021 LA & AASIST & $\sim$0.02 & SOTA at publication \\
  \rowcolor{altrow}
  Tak et al.\ (2022) &
  ASVspoof 2021 LA & RawBoost + ResNet & $\sim$0.015 & Data augmentation \\
  \bottomrule
\end{tabular}
\captionof{table}{Indicative EER of key published anti-spoofing systems. Values are
approximate; exact results vary with evaluation protocol. This table represents the
published state of the art at the time of the Phase 1 literature survey.}
\label{tab:lit_comparison}
\end{center}

\subsection{Summary of Findings}

\begin{keybox}{Key Takeaways --- Deep Learning}
  \begin{itemize}
    \item CNN-based systems match classical CQCC baselines but are limited by global
      temporal pooling, which discards prosodic dynamics.
    \item CNN-LSTM hybrids with frequency-only pooling consistently outperform CNNs
      alone by modelling temporal evolution of spectral features.
    \item Self-supervised models (wav2vec 2.0) and graph attention models (AASIST)
      represent the current frontier, achieving EER below 0.02 on ASVspoof 2021.
    \item The field has progressed rapidly from EER $\approx 0.10$ (CQCC+GMM, 2019)
      to EER $\approx 0.015$ (RawBoost+ResNet, 2022) on standardised benchmarks.
  \end{itemize}
\end{keybox}

% ═════════════════════════════════════════════════════════════════════════════
\section{Robustness, Codec Effects, and Domain Mismatch}
% ═════════════════════════════════════════════════════════════════════════════

\subsection{Channel Mismatch and Codec Compression}

A critical finding across the ASVspoof challenge series is that systems trained on clean
benchmark data degrade significantly under realistic channel conditions:

\begin{itemize}
  \item Yi et al.\ (2022), through the ADD 2022 challenge, explicitly evaluated systems
    under telephone channel conditions, finding EER increases of 5--15 percentage points
    compared to clean evaluation.
  \item G.711 $\mu$-law codec compression (8\,kHz, used in PSTN and banking IVR systems)
    introduces quantisation noise and spectral smoothing that modifies both genuine speech
    characteristics and the synthetic artefacts that detectors rely upon.
  \item ADD 2022 found that codec compression causes moderate and predictable degradation,
    while additive noise (particularly at low SNR) causes more severe and less predictable
    degradation. The consensus finding is that \textbf{noise is a more serious deployment
    threat than codec compression} for anti-spoofing systems in telephony environments.
\end{itemize}

\subsection{Cross-System Generalisation}

Kinnunen et al.\ (2020) investigated how well systems trained on ASVspoof 2019 TTS
systems generalise to newer or unseen TTS architectures:

\begin{itemize}
  \item Systems achieving near-zero EER on ASVspoof 2019 test systems sometimes failed
    on newer TTS architectures that were not present in the training data.
  \item This generalisation failure arises when models learn \textbf{system-specific
    artefacts} (peculiar to particular TTS implementations) rather than \textbf{general
    properties} common to all neural TTS vocoders.
  \item Models trained on \textbf{diverse} sets of TTS systems, covering multiple
    synthesis paradigms (statistical parametric, attention-based, flow-based, GAN-based),
    tend to generalise better to unseen systems than those trained on a narrow set.
\end{itemize}

\subsection{Data Augmentation Strategies}

Tak et al.\ (2022) demonstrated that \textbf{RawBoost} --- a signal-level data
augmentation technique applied during training --- significantly improves robustness:

\begin{itemize}
  \item RawBoost applies random combinations of coloured noise, impulsive noise, and
    stationary background noise to training waveforms, forcing the model to learn
    features that are stable under noise variation.
  \item This training-time augmentation approach outperformed the application of noise
    reduction at inference time, showing that robustness must be built into the model
    rather than compensated for at deployment.
  \item RawBoost combined with ResNet achieved EER $\approx 0.015$ on ASVspoof 2021 LA
    --- the best published result at the time of this survey.
\end{itemize}

\subsection{Summary of Findings}

\begin{keybox}{Key Takeaways --- Robustness and Domain Mismatch}
  \begin{itemize}
    \item Training on clean data and hoping for channel generalisation does not work ---
      EER increases of 5--15 percentage points are reported under telephone conditions.
    \item Noise is a more severe deployment threat than codec compression alone.
      Both must be addressed, but noise deserves priority attention.
    \item Cross-system generalisation is a known failure mode. Training on a diverse set
      of TTS systems covering multiple synthesis paradigms improves robustness to
      unseen systems.
    \item Training-time noise augmentation is more effective than inference-time noise
      reduction for building robust detectors.
  \end{itemize}
\end{keybox}

% ═════════════════════════════════════════════════════════════════════════════
\section{Gaps in the Literature and Project Positioning}
% ═════════════════════════════════════════════════════════════════════════════

\subsection{Identified Gaps}

Despite the maturity of the ASVspoof benchmark and a large body of published systems,
three gaps in the existing literature are directly relevant to this project:

\begin{enumerate}
  \item \textbf{Banking telephony simulation:} The overwhelming majority of published
    work evaluates on clean studio-quality recordings, or on the specific ASVspoof 2021
    evaluation channel. There is comparatively little published work that explicitly
    simulates banking IVR telephony conditions: 8\,kHz sampling rate, G.711 $\mu$-law
    codec, and short transactional utterances (3--8 seconds of banking phrases rather
    than long read sentences).

  \item \textbf{Controlled head-to-head comparison:} Most published work focuses
    exclusively on either classical feature-based systems or deep learning systems.
    Few studies directly compare both under \emph{identical data and conditions}, making
    it difficult to quantify the practical benefit of deep learning over well-engineered
    handcrafted features in a specific deployment context.

  \item \textbf{Robustness degradation curves:} Published results typically report a
    single performance number at a specific evaluation condition. Reporting \emph{continuous
    degradation curves} as a function of SNR level and codec type is rare in the
    anti-spoofing literature, yet this information is essential for deployment engineering
    decisions (e.g., setting operating thresholds under expected noise conditions).
\end{enumerate}

\subsection{Design Decisions Motivated by the Literature}

Table~\ref{tab:decisions} maps each design decision in this project's methodology to
the specific literature finding that motivated it. These decisions govern the dataset
construction, feature engineering, and model development planned for Phases 2 and 3.

\medskip
\begin{center}
\begin{tabular}{L{3.8cm}L{5.5cm}L{4.7cm}}
  \toprule
  \rowcolor{headerbg}
  \color{white}\textbf{Design Decision} &
  \color{white}\textbf{Literature Motivation} &
  \color{white}\textbf{How It Will Be Implemented} \\
  \midrule
  40 MFCC coefficients (not 13) &
  Sahidullah et al.\ (2015): fine spectral structure discriminates vocoders &
  \texttt{n\_mfcc=40} in \texttt{config.yaml}; 257-dim vector \\
  \rowcolor{altrow}
  Delta + delta-delta MFCCs &
  Sahidullah et al.\ (2015): temporal dynamics are the primary discriminative signal &
  120 of 257 dimensions encode temporal rate-of-change \\
  Phase difference features &
  Patel \& Patil (2015): phase coherence artefacts invisible in magnitude spectrum &
  2 dimensions; supplements MFCC information \\
  \rowcolor{altrow}
  8\,kHz + G.711 preprocessing &
  ASVspoof 2021; Yi et al.\ (2022): channel mismatch degrades clean-trained models &
  All utterances processed through codec simulation pipeline \\
  EER as primary metric &
  ASVspoof series standard; balances FAR and FRR at single threshold &
  Threshold set at EER point on validation set \\
  \rowcolor{altrow}
  CNN-LSTM as planned DL architecture &
  Multiple studies: temporal modelling via LSTM improves over CNN-only &
  2-layer BiLSTM on time-preserving spectrogram front-end \\
  Edge-TTS as held-out test &
  Kinnunen et al.\ (2020): cross-system generalisation is a known failure mode &
  1{,}133 clips from 20 commercial voices; zero training overlap \\
  \rowcolor{altrow}
  5/10/15/20\,dB noise variants &
  Yi et al.\ (2022); Tak et al.\ (2022): noise is primary deployment threat &
  Dataset pipeline produces augmented variants; planned for Phase 2 training \\
  Diverse TTS training set &
  Kinnunen et al.\ (2020): diverse training improves generalisation &
  ASVspoof 2019 (6 systems) + VITS + Tacotron2 \\
  \rowcolor{altrow}
  SVM as classical baseline &
  ASVspoof 2019 official baseline; well-validated in the literature &
  RBF kernel with EER-optimal threshold on validation set \\
  \bottomrule
\end{tabular}
\captionof{table}{Literature-to-implementation design decision mapping. These decisions
govern the methodology for Phases 2 and 3.}
\label{tab:decisions}
\end{center}

\begin{notebox}
  \textbf{Note on scope:} The decisions in Table~\ref{tab:decisions} describe
  \emph{planned} methodology for Phase 2, informed by the Phase 1 literature survey.
  They are included here to make the traceability between literature findings and
  project design explicit. The actual implementation and results will be reported in
  the Phase 2 interim report.
\end{notebox}

% ═════════════════════════════════════════════════════════════════════════════
\section{Summary}
% ═════════════════════════════════════════════════════════════════════════════

This literature review has surveyed the field of synthetic speech detection from the
foundational ASVspoof 2015--2021 challenges through to the current state of the art.
The principal conclusions for this project are:

\begin{enumerate}
  \item \textbf{Feature design:} Delta-MFCC coefficients (40 coefficients, with
    $\Delta$ and $\Delta\Delta$) are the strongest discriminators between genuine and
    synthetic speech. Phase-based features provide complementary, non-redundant
    information. Both will be included in the planned handcrafted feature vector.

  \item \textbf{Classifier choice:} SVM with RBF kernel is the appropriate classical
    baseline, consistent with ASVspoof field convention. CNN-LSTM hybrids are the
    preferred deep learning architecture, motivated by consistent superiority over
    CNN-only systems in the temporal modelling of spoofing artefacts.

  \item \textbf{Telephony simulation:} Training under realistic banking channel
    conditions (8\,kHz, G.711) from the outset is essential, not optional. The
    literature clearly shows that clean-trained models fail under channel mismatch.

  \item \textbf{Robustness evaluation:} Noise is the primary deployment threat.
    Degradation curves across SNR levels (5--20\,dB) will be produced in Phase 2 to
    quantify this empirically for the five planned models.

  \item \textbf{Generalisation test:} A held-out system (Microsoft Edge-TTS) not
    present in any published anti-spoofing training set will be used to measure
    cross-system generalisation directly, addressing the gap identified by Kinnunen
    et al.\ (2020).
\end{enumerate}

% ═════════════════════════════════════════════════════════════════════════════
\section*{References}
\addcontentsline{toc}{section}{References}
% ═════════════════════════════════════════════════════════════════════════════

\begin{itemize}[leftmargin=2em, itemsep=5pt]
  \item Alzantot, M., Wang, Z., and Srivastava, M.B. (2019). Deep Residual Neural
    Networks for Audio Spoofing Detection.
    \textit{Proc.\ Interspeech 2019}, 1078--1082.

  \item Baevski, A., Zhou, Y., Mohamed, A., and Auli, M. (2020). wav2vec 2.0: A
    Framework for Self-Supervised Learning of Speech Representations.
    \textit{Advances in Neural Information Processing Systems (NeurIPS)}, 33.

  \item Kim, J., Kong, J., and Son, J. (2021). Conditional Variational Autoencoder
    with Adversarial Learning for End-to-End Text-to-Speech. \textit{ICML 2021}.

  \item Kinnunen, T., Sahidullah, M., Falcone, M., et al. (2020). t-DCF: a Detection
    Cost Function for the Tandem Assessment of Spoofing Countermeasures and Automatic
    Speaker Verification.
    \textit{IEEE/ACM Transactions on Audio, Speech, and Language Processing}.

  \item McFee, B., Raffel, C., Liang, D., et al. (2015). librosa: Audio and Music
    Signal Analysis in Python.
    \textit{Proceedings of the 14th Python in Science Conference}, 18--25.

  \item Nautsch, A., Wang, X., Evans, N., et al. (2021). ASVspoof 2019: Spoofing
    Countermeasures for the Detection of Synthesized, Converted and Replayed Speech.
    \textit{IEEE Transactions on Biometrics, Behavior, and Identity Science}.

  \item Panayotov, V., Chen, G., Povey, D., and Khudanpur, S. (2015). Librispeech:
    An ASR Corpus Based on Public Domain Audio Books. \textit{ICASSP 2015}.

  \item Patel, T.B. and Patil, H.A. (2015). Combining Evidences from Mel Cepstral,
    Cochlear Filter Cepstral and Instantaneous Frequency Features for Detection of
    Natural vs.\ Spoofed Speech. \textit{Interspeech 2015}, 2062--2066.

  \item Sahidullah, M., Kinnunen, T., and Hanilci, C. (2015). A Comparison of Features
    for Synthetic Speech Detection.
    \textit{Interspeech 2015}, 2087--2091.

  \item Tak, H., Patino, J., Todisco, M., et al. (2022). End-to-End Spectro-Temporal
    Graph Attention Networks for Speaker Verification Anti-Spoofing and Speech Deepfake
    Detection. \textit{ASVspoof 2021 Workshop}.

  \item Todisco, M., Wang, X., Vestman, V., et al. (2019). ASVspoof 2019: Future
    Horizons in Spoofed and Fake Speech Detection. \textit{Interspeech 2019}.

  \item van den Oord, A., Dieleman, S., Zen, H., et al. (2016). WaveNet: A Generative
    Model for Raw Audio. \textit{arXiv:1609.03499}.

  \item Wang, T., Yi, J., Fu, R., et al. (2021). AASIST: Audio Anti-Spoofing using
    Integrated Spectro-Temporal Graph Attention Networks. \textit{ICASSP 2022}.

  \item Wang, Y., Skerry-Ryan, R.J., Stanton, D., et al. (2017). Tacotron: Towards
    End-to-End Speech Synthesis. \textit{Interspeech 2017}.

  \item Yi, J., Fu, R., Tao, J., et al. (2022). ADD 2022: The First Audio Deep
    Synthesis Detection Challenge. \textit{ICASSP 2022}.
\end{itemize}

\end{document} 